\section{Demand, disruption and recovery follow different paths}
\label{sec:apple-disruption}

An open supply route is useful only if a customer still wants the product.
Conversely, a customer may be ready to buy while a supplier cannot deliver a
critical component. The operating scenario introduced above must therefore
describe both demand and the ability to serve it. Apple's annual risk discussion
connects global and regional economic conditions to consumer spending, the
financial condition of its business partners and its own access to finance.
Its separate interruption discussion describes obstacles to manufacture,
delivery and service provision \citep[Risk Factors, printed pp.~5--6]{apple2024}.
These passages give reasons to retain several possible operating histories
behind the same reported financial totals.

Apple says that sales outside the United States account for a majority of net
sales and that a majority of its supplier facilities are outside the United
States. The two statements concern different populations. A sales region does
not locate a factory, and a factory's location does not identify the customers
whose orders depend on it. The segment definitions examined earlier generally
attribute sales by customer and retail-store location, while manufacturing
purchase obligations describe purchases from outsourcing partners and component
suppliers \citep[Note 13, printed p.~46; Liquidity and Capital Resources,
printed p.~25]{apple2024}. Linking the two requires the actual product,
component, supplier and route relationships. Geographic labels alone cannot
create them.

The macroeconomic passage describes several possible effects of the same
economic conditions. Recession, unemployment, inflation, tighter credit,
higher interest rates and currency fluctuations can weaken consumer confidence
and spending. Fiscal and monetary policy, market volatility and changes in
income or asset values can also affect demand. Suppliers, manufacturers,
logistics providers, distributors, carriers, other channel partners and
developers may instead face financial instability, difficulty obtaining credit
or insolvency. Apple further identifies trade-receivable collectibility,
derivative-counterparty and financial-institution failure, limited access to
new debt, reduced liquidity and declining financial-instrument fair values.
The disclosure permits a scenario to connect these channels through common
conditions; it supplies neither a joint probability law nor a requirement
that every channel deteriorate together.

An interruption raises a different operating question: which activity becomes
unavailable, and for how long? Apple's discussion includes political and trade
disputes, geopolitical conflict, terrorism, natural disasters, public-health
issues and industrial accidents. It says substantially all manufacturing is
performed wholly or partly by outsourcing partners located primarily in China
mainland, India, Japan, South Korea, Taiwan and Vietnam. Trade restrictions can
raise input costs, require different suppliers or business relationships, or
lead the company to cease offering affected products, services or third-party
applications. Changing the supply chain may itself be costly and disruptive,
and restrictions may arrive with little notice. Apple reports that past
U.S.--China tensions led to tariffs and other restrictions affecting its
business. This historical assertion is distinct from its conditional warning
about further escalation; a future scenario still needs a specified measure,
effective date and affected trade.

The same passage describes facilities in disaster-prone locations and asserts
that climate change is increasing the frequency or intensity of certain
natural disasters and extreme weather. It also identifies fire, power
shortages, nuclear and other industrial accidents, hostile acts, ransomware,
labor disputes and health emergencies. These are Apple's stated exposures,
without facility-specific hazard estimates. Its requirement that suppliers
maintain safe workplaces is separate from whether an accident occurs. The
company reports that major public-health issues, including COVID-19, have
adversely affected demand, operations and distribution, and warns of possible
future effects. Travel restrictions, freight limitations and interruptions to
production ramps supply concrete mechanisms; the paragraph does not quantify
an historical loss or a future duration.

Consider two explicitly assumed histories for a particular product. In one,
the supply route remains available but customers reduce orders. In another,
orders remain but a critical component cannot arrive. The first requires a
demand response; the second requires an inventory and delivery response.
They can also occur together. To compare the histories, the model needs
product-level orders, available components and finished goods, lead times,
contractual commitments and feasible management actions. An assumed alternate
supplier must have capacity and suitable terms before it becomes an available
response. A sales delay or lost order follows from those facts and assumptions;
it cannot be obtained by applying one percentage to every reported balance.

Recovery adds another date. Apple warns that an interruption can require
substantial time and expenditure to resume operations, with lost sales and
greater consequences where critical components have single or limited sources.
Note 12 separately explains that even components available from several
suppliers can face industry-wide shortages and commodity-price fluctuations
\citep[Note 12, printed p.~46]{apple2024}. Ending an assumed interruption
therefore need not restore production or sales immediately. Recovery spending
can precede or follow resumed deliveries. The model must retain the dates of
the interruption, restored capacity, deliveries, expenditure and settlement
separately; the applicable accounting treatment then determines the journals.
The insurance qualification already examined in the litigation discussion
does not collapse those dates into a net loss or guarantee reimbursement.

The reported receivables, investments, debt and liquidity resources constrain
particular dated histories, but they do not determine these operating
responses. The investment sensitivity discussed above measures a specified
interest-rate effect. Chapter~\ref{ch:financing-groups} examines the distinct
debt and currency measures; none identifies the broader disruption response.
The interim's incorporation of annual risk factors, except for the changes it
sets out, likewise does not establish unchanged facilities, orders or funding
conditions \citep[Part II, Item 1A, printed p.~19]{apple2025q1}.
For the historical record, the useful distinction is now precise: retain
Apple's assertions, the relationships needed to interpret them, and the
analyst's scenario assumptions as separate information. Alternative demand,
supply and recovery histories can then remain visible until evidence narrows
them.
